% Options for packages loaded elsewhere
% Options for packages loaded elsewhere
\PassOptionsToPackage{unicode}{hyperref}
\PassOptionsToPackage{hyphens}{url}
\PassOptionsToPackage{dvipsnames,svgnames,x11names}{xcolor}
%
\documentclass[
  letterpaper,
  DIV=11,
  numbers=noendperiod]{scrartcl}
\usepackage{xcolor}
\usepackage{amsmath,amssymb}
\setcounter{secnumdepth}{-\maxdimen} % remove section numbering
\usepackage{iftex}
\ifPDFTeX
  \usepackage[T1]{fontenc}
  \usepackage[utf8]{inputenc}
  \usepackage{textcomp} % provide euro and other symbols
\else % if luatex or xetex
  \usepackage{unicode-math} % this also loads fontspec
  \defaultfontfeatures{Scale=MatchLowercase}
  \defaultfontfeatures[\rmfamily]{Ligatures=TeX,Scale=1}
\fi
\usepackage{lmodern}
\ifPDFTeX\else
  % xetex/luatex font selection
\fi
% Use upquote if available, for straight quotes in verbatim environments
\IfFileExists{upquote.sty}{\usepackage{upquote}}{}
\IfFileExists{microtype.sty}{% use microtype if available
  \usepackage[]{microtype}
  \UseMicrotypeSet[protrusion]{basicmath} % disable protrusion for tt fonts
}{}
\makeatletter
\@ifundefined{KOMAClassName}{% if non-KOMA class
  \IfFileExists{parskip.sty}{%
    \usepackage{parskip}
  }{% else
    \setlength{\parindent}{0pt}
    \setlength{\parskip}{6pt plus 2pt minus 1pt}}
}{% if KOMA class
  \KOMAoptions{parskip=half}}
\makeatother
% Make \paragraph and \subparagraph free-standing
\makeatletter
\ifx\paragraph\undefined\else
  \let\oldparagraph\paragraph
  \renewcommand{\paragraph}{
    \@ifstar
      \xxxParagraphStar
      \xxxParagraphNoStar
  }
  \newcommand{\xxxParagraphStar}[1]{\oldparagraph*{#1}\mbox{}}
  \newcommand{\xxxParagraphNoStar}[1]{\oldparagraph{#1}\mbox{}}
\fi
\ifx\subparagraph\undefined\else
  \let\oldsubparagraph\subparagraph
  \renewcommand{\subparagraph}{
    \@ifstar
      \xxxSubParagraphStar
      \xxxSubParagraphNoStar
  }
  \newcommand{\xxxSubParagraphStar}[1]{\oldsubparagraph*{#1}\mbox{}}
  \newcommand{\xxxSubParagraphNoStar}[1]{\oldsubparagraph{#1}\mbox{}}
\fi
\makeatother


\usepackage{longtable,booktabs,array}
\newcounter{none} % for unnumbered tables
\usepackage{calc} % for calculating minipage widths
% Correct order of tables after \paragraph or \subparagraph
\usepackage{etoolbox}
\makeatletter
\patchcmd\longtable{\par}{\if@noskipsec\mbox{}\fi\par}{}{}
\makeatother
% Allow footnotes in longtable head/foot
\IfFileExists{footnotehyper.sty}{\usepackage{footnotehyper}}{\usepackage{footnote}}
\makesavenoteenv{longtable}
\usepackage{graphicx}
\makeatletter
\newsavebox\pandoc@box
\newcommand*\pandocbounded[1]{% scales image to fit in text height/width
  \sbox\pandoc@box{#1}%
  \Gscale@div\@tempa{\textheight}{\dimexpr\ht\pandoc@box+\dp\pandoc@box\relax}%
  \Gscale@div\@tempb{\linewidth}{\wd\pandoc@box}%
  \ifdim\@tempb\p@<\@tempa\p@\let\@tempa\@tempb\fi% select the smaller of both
  \ifdim\@tempa\p@<\p@\scalebox{\@tempa}{\usebox\pandoc@box}%
  \else\usebox{\pandoc@box}%
  \fi%
}
% Set default figure placement to htbp
\def\fps@figure{htbp}
\makeatother





\setlength{\emergencystretch}{3em} % prevent overfull lines

\providecommand{\tightlist}{%
  \setlength{\itemsep}{0pt}\setlength{\parskip}{0pt}}



 


\KOMAoption{captions}{tableheading}
\makeatletter
\@ifpackageloaded{caption}{}{\usepackage{caption}}
\AtBeginDocument{%
\ifdefined\contentsname
  \renewcommand*\contentsname{Table of contents}
\else
  \newcommand\contentsname{Table of contents}
\fi
\ifdefined\listfigurename
  \renewcommand*\listfigurename{List of Figures}
\else
  \newcommand\listfigurename{List of Figures}
\fi
\ifdefined\listtablename
  \renewcommand*\listtablename{List of Tables}
\else
  \newcommand\listtablename{List of Tables}
\fi
\ifdefined\figurename
  \renewcommand*\figurename{Figure}
\else
  \newcommand\figurename{Figure}
\fi
\ifdefined\tablename
  \renewcommand*\tablename{Table}
\else
  \newcommand\tablename{Table}
\fi
}
\@ifpackageloaded{float}{}{\usepackage{float}}
\floatstyle{ruled}
\@ifundefined{c@chapter}{\newfloat{codelisting}{h}{lop}}{\newfloat{codelisting}{h}{lop}[chapter]}
\floatname{codelisting}{Listing}
\newcommand*\listoflistings{\listof{codelisting}{List of Listings}}
\makeatother
\makeatletter
\makeatother
\makeatletter
\@ifpackageloaded{caption}{}{\usepackage{caption}}
\@ifpackageloaded{subcaption}{}{\usepackage{subcaption}}
\makeatother
\makeatletter
\@ifpackageloaded{fontawesome5}{}{\usepackage{fontawesome5}}
\makeatother
\usepackage{bookmark}
\IfFileExists{xurl.sty}{\usepackage{xurl}}{} % add URL line breaks if available
\urlstyle{same}
\hypersetup{
  pdftitle={STA 198/GLHLTH 198 - Introduction to Global Health Data Science},
  colorlinks=true,
  linkcolor={blue},
  filecolor={Maroon},
  citecolor={Blue},
  urlcolor={Blue},
  pdfcreator={LaTeX via pandoc}}


\title{STA 198/GLHLTH 198 - Introduction to Global Health Data Science}
\usepackage{etoolbox}
\makeatletter
\providecommand{\subtitle}[1]{% add subtitle to \maketitle
  \apptocmd{\@title}{\par {\large #1 \par}}{}{}
}
\makeatother
\subtitle{Fall 2026}
\author{}
\date{}
\begin{document}
\maketitle

\renewcommand*\contentsname{Table of contents}
{
\setcounter{tocdepth}{3}
\tableofcontents
}

\section{Short Summary}\label{short-summary}

\subsection{Course Overview}\label{course-overview}

Learn to explore, visualize, and analyze data using R and RStudio
through real-world global health problems. Study best practices in data
wrangling, visualization, modeling, and statistical inference.
Communicate results in plan language accessible to non-experts.
Distinguish good evidence from weak or misleading evidence using
statistical reasoning.

\subsection{Required Materials}\label{required-materials}

\begin{itemize}
\tightlist
\item
  \textbf{Textbooks} (free online): \emph{R for Data Science, 2e} and
  \emph{Introduction to Modern Statistics}; (free in library reserves
  and optional): \emph{Principles of Biostatistics}
\item
  \textbf{Computing}: Laptop required for all classes; R accessed
  through Duke containers, no installation needed
\item
  \textbf{Platform}: Course website at
  \href{https://sta198-f26-2.github.io}{sta198-f26-2.github.io} --
  everything you need is there or linked from there!
\end{itemize}

\subsection{Assessment \& Grading}\label{assessment-grading}

{\def\LTcaptype{none} % do not increment counter
\begin{longtable}[]{@{}
  >{\raggedright\arraybackslash}p{(\linewidth - 4\tabcolsep) * \real{0.3889}}
  >{\raggedright\arraybackslash}p{(\linewidth - 4\tabcolsep) * \real{0.2917}}
  >{\raggedright\arraybackslash}p{(\linewidth - 4\tabcolsep) * \real{0.3194}}@{}}
\toprule\noalign{}
\begin{minipage}[b]{\linewidth}\raggedright
Component
\end{minipage} & \begin{minipage}[b]{\linewidth}\raggedright
Weight
\end{minipage} & \begin{minipage}[b]{\linewidth}\raggedright
Details
\end{minipage} \\
\midrule\noalign{}
\endhead
\bottomrule\noalign{}
\endlastfoot
Lectures & 5\% & Attendance/participation (20 non-exam lectures minimum
for full score) \\
Labs & 10\% & Individual or team-based, completed in class (lowest 2
dropped) \\
Exam 1 & 25\% & In-class \\
Exam 2 & 25\% & In-class \\
Exam 3 & 25\% & In-class \\
Final & Replace lowest exam if desired & In-class \\
\end{longtable}
}

\subsection{Key Policies}\label{key-policies}

\subsubsection{Academic Honesty}\label{academic-honesty}

\begin{itemize}
\tightlist
\item
  \textbf{Individual work}: Exams must be completed alone
\item
  \textbf{Team collaboration}: Expected for some labs
\item
  \textbf{AI tools}: Allowed for code assistance with proper citation
\item
  \textbf{Online resources}: Permitted with explicit citation
\end{itemize}

\subsubsection{Deadlines \& Late Work}\label{deadlines-late-work}

\begin{itemize}
\tightlist
\item
  \textbf{Labs}: No late work (completed in class)
\item
  \textbf{Exams}: No make-ups (final can substitute for lowest score)
\end{itemize}

\subsubsection{Attendance}\label{attendance}

\begin{itemize}
\tightlist
\item
  \textbf{Lectures}: Mandatory attendance
\item
  \textbf{Labs}: Mandatory attendance
\end{itemize}

\subsection{Important Dates}\label{important-dates}

\begin{itemize}
\tightlist
\item
  \textbf{Classes begin}: Aug 24
\item
  \textbf{Exam 1}: Sept 30
\item
  \textbf{Exam 2}: Nov 4
\item
  \textbf{Exam 3}: Dec 2
\item
  \textbf{Classes end}: Dec 4
\item
  \textbf{Final exam}: Dec 11
\end{itemize}

\subsection{Getting Help}\label{getting-help}

\begin{itemize}
\tightlist
\item
  \textbf{Office hours}: Regular TA and instructor availability
\item
  \textbf{Ed Discussion}: Online forum for course questions
\item
  \textbf{Email}: Include ``STA 198'' in subject line
\item
  \textbf{Accommodations}: Contact SDAO for academic accommodations
\end{itemize}

\section{Details}\label{details}

\subsubsection{Textbooks}\label{textbooks}

The first two books are \textbf{freely available online}.

\begin{itemize}
\item
  \href{https://r4ds.hadley.nz/}{R for Data Science, 2e}, Wickham,
  Çetinkaya-Rundel, Grolemund. O'Reilly, 2nd edition, 2023. Hard copy
  available
  \href{https://www.amazon.com/dp/1492097403?&tag=hadlwick-20}{on
  Amazon}.
\item
  \href{https://openintro-ims.netlify.app/}{Introduction to Modern
  Statistics}, Çetinkaya-Rundel, Hardin. OpenIntro Inc., 2nd Edition,
  2023. Hard copy available
  \href{https://openintro.org/go?id=ims2_bw_pb&referrer=/book/ims/index.php}{on
  Amazon}.
\end{itemize}

The last book is not required. Three copies are on reserve at the
library, and the problem sets and ``further applications'' worked
examples in each chapter are \emph{quite} useful for students wanting
more examples and problems (I suggest everyone check out the further
applications!).

\begin{itemize}
\tightlist
\item
  \href{https://openintro-ims.netlify.app/}{Principles of
  Biostatistics}, Pagano, Gauvreau, and Mattie. Taylor and Francis, 3rd
  Edition, 2022.
\end{itemize}

\subsubsection{Computing}\label{computing}

You will need a laptop you can bring to lecture and lab for this course.
We will use the statistical software R. Students will be able to access
R through Docker containers provided by Duke Office of Information
Technology. See the \href{computing-access.qmd}{computing page} for more
information.

\subsubsection{Accessibility}\label{accessibility}

If any portion of the course is not accessible to you due to challenges
with technology or the course format, please let me know so we can make
appropriate accommodations.

The \href{https://access.duke.edu/students}{Student Disability Access
Office (SDAO)} is available to ensure that students can engage with
their courses and related assignments. Students should contact the SDAO
to \href{https://access.duke.edu/requests}{request or update
accommodations} under these circumstances. If you have an exam
accommodation, please schedule a seating there for the three exams and
final as soon as possible.

\subsubsection{Communication}\label{communication}

All lecture notes, assignment instructions, an up-to-date schedule, and
other course materials may be found on the course website:
\href{https://sta198-f26-2.github.io/}{sta198-f26-2.github.io}.

Announcements will periodically be emailed through Canvas Announcements.
Please check your email regularly to ensure you have the latest
announcements for the course.

\subsubsection{Email}\label{email}

If you have questions about assignment extensions, accommodations, or
any other matter not appropriate for the class discussion forum, please
email me directly at
\href{mailto:ahh24@duke.edu}{\nolinkurl{ahh24@duke.edu}}. \textbf{If you
do so, please include ``STA 198'' in the subject line.} Barring
extenuating circumstances, I will respond to STA 198 emails within 48
hours, Monday through Friday. Response time may be slower for emails
sent Friday evening through Sunday.

\subsection{Tips for success}\label{tips-for-success}

Your success in this course depends very much on you and the effort you
put into it. Your TA and I will help you by providing you with materials
answering questions, and setting a pace, but for this to work, you must
do the following:

\begin{enumerate}
\def\labelenumi{\arabic{enumi}.}
\tightlist
\item
  \textbf{Keep up with the material.}
\end{enumerate}

Come prepared so you can deeply engage with the material during lectures
and labs. The schedule lists extensive additional resources if you need
additional support -- take advantage so you don't fall behind.

\begin{enumerate}
\def\labelenumi{\arabic{enumi}.}
\setcounter{enumi}{1}
\tightlist
\item
  \textbf{Be present and engaged in every lecture and lab.}
\end{enumerate}

In lectures ask questions and participate in discussions. In labs, work
on the lab exercises, ask questions, and collaborate with your
teammates. If you miss a class, make sure to catch up on the material
before the next class.

\begin{enumerate}
\def\labelenumi{\arabic{enumi}.}
\setcounter{enumi}{2}
\tightlist
\item
  \textbf{Ask questions.}
\end{enumerate}

As often as you can. In class, out of class. Ask me, ask the TAs, ask
your friends, ask the person sitting next to you. This will help you
more than anything else. If you get a question wrong on an assessment,
ask us why. If you're not sure about the lab, ask. If you hear something
on the news that sounds related to what we discussed, ask. If the
reading is confusing, ask.

\begin{enumerate}
\def\labelenumi{\arabic{enumi}.}
\setcounter{enumi}{3}
\tightlist
\item
  \textbf{Do the work.}
\end{enumerate}

The earlier you start, the better. It's not enough to just mechanically
plow through material. You should ask yourself how things relate to
earlier material and imagine how they might be changed (to make
questions for an exam, for example).

\begin{enumerate}
\def\labelenumi{\arabic{enumi}.}
\setcounter{enumi}{4}
\tightlist
\item
  \textbf{Don't procrastinate.}
\end{enumerate}

The content builds upon what was taught in previous weeks, so if
something is confusing to you in Week 2, Week 3 will become more
confusing, Week 4 even worse, etc. Don't let the week end with
unanswered questions. But if you find yourself falling behind and not
knowing where to begin asking, come to office hours and work with a
member of the teaching team to help you identify a good (re)starting
point.

\subsection{Getting help}\label{getting-help-1}

\begin{itemize}
\tightlist
\item
  If you have a question during the lecture or lab, feel free to ask it!
  There are likely other students with the same question, so by asking,
  you will create a learning opportunity for everyone.
\item
  The teaching team is here to help you be successful in the course. You
  are encouraged to attend office hours to ask questions about the
  course content and assignments. Many questions are most effectively
  answered as you discuss them with others, so office hours are a
  valuable resource. Please use them!
\item
  Outside of class and office hours, any general questions about course
  content or assignments should be posted on the class discussion forum,
  \href{https://edstem.org/}{Ed Discussion}. There is a chance another
  student has already asked a similar question, so please check the
  other posts on the forum before adding a new question. If you know the
  answer to a question, I encourage you to respond!
\end{itemize}

Check out the \href{course-support.html}{Support} tab for more
resources.

\subsection{Course components}\label{course-components}

\subsubsection{Lectures}\label{lectures}

Lectures designed to be interactive, so you gain experience applying new
concepts and learning from each other. My role as instructor is to
introduce you to new methods, tools, and techniques, but it is up to you
to take them and use them. A lot of what you do in this course will
involve writing code, and coding is a skill that is best learned by
doing. Therefore, many lectures will feature application exercises that
will serve as opportunities to practice what you're learning as you're
learning it and be great preparation for the assignments and exams.

You are expected to bring a laptop (or Chromebook) to each class so that
you can participate in the application activities. Please ensure your
device is fully charged before you come to class, as the number of
outlets in the classroom will not be sufficient to accommodate everyone.
A tablet also works, but the user experience will be much smoother on a
laptop.

\subsubsection{Labs}\label{labs}

Labs are designed to be hands-on at all times, therefore you're
similarly expected to bring a laptop to each lab session.

During labs you will get a brief introduction to the week's assignments
from your TA and then work on the lab exercise, either individually or
with your teammates.

\subsubsection{Teams}\label{teams}

Many labs will involve teamwork. All team members are expected to
contribute equally to the completion of the lab, and you will be asked
to evaluate your team members throughout the semester. Failure to
adequately contribute to any component will result in a penalty to your
mark relative to the team's overall mark. You are expected to use the
provided GitHub repository as the central collaborative platform.
Commits to this repository will be used as a metric (one of several) of
each team member's relative contribution to each lab.

\subsubsection{Activities \& Assessment}\label{activities-assessment}

\paragraph{Lectures}\label{lectures-1}

Attendance and participation will be monitored throughout the semester.
Your attendance and participation score will be calculated as the
percentage of lectures in which you engaged in person (i.e., sleeping or
online shopping does not count, and you must participate in class
activities). The denominator used for this calculation will be 22 (we
have 22 non-exam lecture days, excluding the first day of class).
Students who are active in 20 or more lectures will receive full credit.
That is, if you actively attend all 22 lectures after the first day of
class and the constellation activity, you will receive 100\%; 20
lectures - 100\%, 19 lectures - \(100*\frac{19}{22}=86\)\%. Note: if you
have an unpleasant communicable disease, please do not share!

\paragraph{Labs}\label{labs-1}

You must be present in lab to complete the lab assignment. There is no
way to make up for missing a lab. You will submit your lab assignments
by pushing your work to your GitHub repository for the lab and
submitting the PDF output on Gradescope \emph{by the end of your lab
session.}

\emph{The lowest two lab grades will be dropped at the end of the
semester, which means you can miss up to two lab sessions with no
penalty.}

\paragraph{Exams}\label{exams}

This course will have three in-class exams. You can demonstrate what
you've learned in the course thus far through these exams. The exams
will focus on both conceptual understanding of the content and
application. The exam's content will be related to the content in
lectures and labs.

\emph{Missed exams: There are no make-up exams. If you miss an exam due
to an illness or similar issue documented with a Dean's Excuse, your
final exam score will replace the missed exam score. If you miss the
Final Exam due to an illness or similar issue documented with a Dean's
Excuse, you will receive an Incomplete and can take the final exam
later.}

\emph{Improvement bonus: For students who take all exams (Exam 1, Exam
2, Exam 3, and the Final Exam), the final exam score will replace the
lowest of the three mid-semester exam scores, if the final exam score is
higher.}

\subsection{Assessment \& Grading}\label{assessment-grading-1}

{\def\LTcaptype{none} % do not increment counter
\begin{longtable}[]{@{}
  >{\raggedright\arraybackslash}p{(\linewidth - 4\tabcolsep) * \real{0.3889}}
  >{\raggedright\arraybackslash}p{(\linewidth - 4\tabcolsep) * \real{0.2917}}
  >{\raggedright\arraybackslash}p{(\linewidth - 4\tabcolsep) * \real{0.3194}}@{}}
\toprule\noalign{}
\begin{minipage}[b]{\linewidth}\raggedright
Component
\end{minipage} & \begin{minipage}[b]{\linewidth}\raggedright
Weight
\end{minipage} & \begin{minipage}[b]{\linewidth}\raggedright
Details
\end{minipage} \\
\midrule\noalign{}
\endhead
\bottomrule\noalign{}
\endlastfoot
Lectures & 5\% & Attendance/participation (20 minimum for full score) \\
Labs & 10\% & Individual or team-based, completed in class (lowest 2
dropped) \\
Exam 1 & 25\% & In-class \\
Exam 2 & 25\% & In-class \\
Exam 3 & 25\% & In-class \\
Final & Replace lowest exam if desired & In-class \\
\end{longtable}
}

The final letter grade will be determined based on the following
thresholds.

{\def\LTcaptype{none} % do not increment counter
\begin{longtable}[]{@{}ll@{}}
\toprule\noalign{}
Letter Grade & Final Course Grade \\
\midrule\noalign{}
\endhead
\bottomrule\noalign{}
\endlastfoot
A & {[}93,100{]} \\
A- & {[}90,93) \\
B+ & {[}87,90) \\
B & {[}83,87) \\
B- & {[}80,83) \\
C+ & {[}77,80) \\
C & {[}73,77) \\
C- & {[}70,73) \\
D+ & {[}67,70) \\
D & {[}63,67) \\
D- & {[}60,63) \\
F & {[}0,60) \\
\end{longtable}
}

\subsection{Course policies}\label{course-policies}

\subsubsection{Duke Community Standard}\label{duke-community-standard}

As a student in this course, you have agreed to uphold the
\href{https://trinity.duke.edu/undergraduate/academic-policies/community-standard-student-conduct}{Duke
Community Standard} and the practices specific to this course.

\subsubsection{Academic honesty}\label{academic-honesty-1}

\textbf{TL;DR: Don't cheat!}

\paragraph{What is allowed and what is
not?}\label{what-is-allowed-and-what-is-not}

Please abide by the following as you work on assignments in this course:

\begin{itemize}
\item
  \textbf{Collaboration:} Only work that is clearly assigned as teamwork
  should be completed collaboratively. On individual assignments, you
  may not directly share work (including code) with another student in
  this class; on team assignments, you may not directly share work
  (including code) with another team. ``Sharing'' includes, but is not
  limited to, messaging, emailing, or otherwise providing your work to
  another student or team.

  \begin{itemize}
  \item
    Labs: Collaboration in teams for team lab assignments is not only
    allowed but expected. You will work together with your lab team to
    complete the lab exercise. However, each student must submit their
    own write-up of the lab exercise, which should reflect their
    understanding and ideas. It's expected that lab submission will be
    similar across team members.
  \item
    Exams: You may not discuss or otherwise work with others on the
    exams while in progress. On exams, collaboration or using
    unauthorized materials will be considered a violation for all
    students involved, regardless of who initiated the sharing.
  \end{itemize}
\item
  \textbf{Use of online resources:} I am well aware that a huge volume
  of code is available on the web to solve any number of problems.
  Unless I explicitly tell you not to use something, the course's policy
  is that you may make use of any online resources, but you must
  explicitly cite where you obtained any code you directly use (or use
  as inspiration). \textbf{Any recycled code that is discovered and is
  not explicitly cited will be treated as plagiarism, resulting in an
  automatic 0 for the relevant portion of the assignment.}
\item
  \textbf{Use of generative artificial intelligence (AI)}: You should
  treat generative AI, such as ChatGPT, like other online resources. Two
  guiding principles govern how to use AI in this course:

  \begin{enumerate}
  \def\labelenumi{\arabic{enumi}.}
  \item
    \emph{Cognitive dimension:} Working with AI should not reduce your
    thinking ability.
  \item
    \emph{Ethical dimension}\textbf{:} Students using AI should be
    transparent about their use and ensure it aligns with academic
    integrity.
  \end{enumerate}

  \begin{itemize}
  \tightlist
  \item
    {\faIcon{check}} \textbf{AI tools for code:} You may use generative
    AI tools when you need help with assignments. However you should
    first attempt to solve the problem yourself. Your submission should
    not be a copy-paste of AI-generated content -- you must edit the
    content to ensure it reflects your understanding, has your voice and
    intellectual input, and conforms with course materials, syntax,
    terminology, and style. Additionally, \textbf{the prompt you use
    should not be copied and pasted directly from the assignment; you
    should create a prompt yourself.}
  \end{itemize}

  You must also explicitly cite work submitted that is based on
  AI-generated content. You may use
  \href{https://guides.lib.monash.edu/c.php?g=219786&p=6972087}{these
  guidelines} to cite AI-generated content. \textbf{The bare minimum
  citation must include the AI tool you're using (e.g., ChatGPT), the
  model the tool uses, the date when you ran the prompt, and a link to
  the full transcript of the session starting with your prompt.} A new
  citation must be included for each exercise where you used AI tools,
  with a new link to the full transceipt of the interaction.
\end{itemize}

While you're welcomed to ask AI tools questions that might help your
learning and understanding in this course, you should be critical of the
answers you receive, as AI-generated content may not always be accurate
or reliable. Use it to supplement your understanding, not as a
substitute for learning. You are ultimately responsible for the work you
turn in; it should reflect \emph{your} understanding of the course
content.

\paragraph{What happens if you violate the academic honesty
policy?}\label{what-happens-if-you-violate-the-academic-honesty-policy}

Any violations in academic honesty standards as outlined in the
\href{https://trinity.duke.edu/undergraduate/academic-policies/community-standard-student-conduct}{Duke
Community Standard} and those specific to this course

\begin{itemize}
\item
  will automatically result in a 0 for the relevant portion or the
  entire assignment or assessment,
\item
  can result in further deductions to your overall course grade (e.g.,
  drop down to the next letter grade or drop down to an F), and
\item
  can be reported to the
  \href{https://students.duke.edu/get-assistance/community-standard/osccs/}{Office
  of Student Conduct \& Community Standards} for further action.
\end{itemize}

Regardless of course delivery format, it is the responsibility of all
students to understand and follow all Duke policies, including academic
integrity (e.g., completing one's own work, following proper citation of
sources, adhering to guidance around group work projects, and more).
Ignoring these requirements is a violation of the Duke Community
Standard. Any questions and/or concerns regarding academic integrity can
be directed to the Office of Student Conduct and Community Standards at
\href{mailto:conduct@duke.edu}{\nolinkurl{conduct@duke.edu}}.

\subsubsection{Regrade requests}\label{regrade-requests}

Regrade requests must be submitted on Gradescope within a week after an
assignment is returned. Regrade requests will be considered if there was
an error in the grade calculation or if a correct answer was mistakenly
marked as incorrect. Requests to dispute the number of points deducted
for an incorrect response will not be considered. Regrade requests are
also not a mechanism for asking for clarification on feedback, those
questions should be brought to office hours. Note that by submitting a
regrade request, the \textbf{entire assignment may be regraded}, which
could potentially result in losing points.

\emph{No grades will be changed after the final exam has been
administered.}

\subsection{Accommodations}\label{accommodations}

\subsubsection{Academic accommodations}\label{academic-accommodations}

If you need accommodations for this class, you will need to register
with the Student Disability Access Office (SDAO) and provide them with
documentation related to your needs. SDAO will work with you to
determine what accommodations are appropriate for your situation. Please
note that accommodations are not retroactive and disability
accommodations cannot be provided until a Faculty Accommodation Letter
has been given to me. Please contact SDAO for more information:
\href{mailto:sdao@duke.edu}{\nolinkurl{sdao@duke.edu}} or
\href{https://access.duke.edu/}{access.duke.edu}. Please note that the
testing center can fill, so you want to go ahead and schedule
accommodations on exam days (overlapping with our exam time) as soon as
possible.

\subsubsection{Religious accommodations}\label{religious-accommodations}

Students are permitted by university policy to be absent from class to
observe a religious holiday. Accordingly, Trinity College of Arts \&
Sciences and the Pratt School of Engineering have established procedures
to be followed by students for notifying their instructors of an absence
necessitated by the observance of a religious holiday. Please submit
requests for religious accommodations at the beginning of the semester
so that we can work to make suitable arrangements well ahead of time.
You can find the policy and relevant notification form here:
\href{https://trinity.duke.edu/undergraduate/academic-policies/religious-holidays}{trinity.duke.edu/undergraduate/academic-policies/religious-holidays}




\end{document}
